%% Q1
Each term (from the third) is the \textbf{sum of all previous terms}: $1+3=4$, $1+3+4=8$, $1+3+4+8=16$, so next $=1+3+4+8+16=\mathbf{32}$.

%% Q2
The differences are $2, 4, 6, 8$, so next $+10$: $30+10=\mathbf{40}$.

%% Q3
Divide successively by 1, 2, 3, 4: $240, 240, 120, 40, \mathbf{10}$.

%% Q4
CENTRAL $\rightarrow$ XVMGIZO uses the \textbf{reverse-alphabet} code (A$\leftrightarrow$Z, B$\leftrightarrow$Y, C$\leftrightarrow$X, \ldots; the two letters add to 27).
PHYSICS: P$\rightarrow$K, H$\rightarrow$S, Y$\rightarrow$B, S$\rightarrow$H, I$\rightarrow$R, C$\rightarrow$X, S$\rightarrow$H, giving \textbf{KSBHRXH}.

%% Q5
Each letter moves one step forward (G$\rightarrow$H, O$\rightarrow$P, \ldots). LIME $\rightarrow$ \textbf{MJNF}.

%% Q6
CHEMISTRY $\rightarrow$ DGFLJRUQZ: the shifts alternate $+1, -1$.
DESIGN: D$+1=$E, E$-1=$D, S$+1=$T, I$-1=$H, G$+1=$H, N$-1=$M, giving \textbf{EDTHHM}.

%% Q7
CARPET $\rightarrow$ XZIKVG is again the reverse-alphabet code. MARKET $\rightarrow$ N, Z, I, P, V, G $=$ \textbf{NZIPVG}.

%% Q8
The quality of research is judged mainly by the \textbf{methodology} followed: the rigour of design, sampling, tools and analysis determines the validity of the findings.

%% Q9
Research is by definition a \emph{systematic} inquiry. So ``research is not always systematic'' is \textbf{not correct}.

%% Q10
The first step is \textbf{identification of the problem}; selection and definition follow.

%% Q11
A good researcher must above all possess an \textbf{inquiring mind} (curiosity).

%% Q12
Per the key, ``quantitative data can be converted into qualitative data'' is treated as \textbf{not correct}. Statement (D), converting qualitative data into quantitative data by coding and counting, is common practice.

%% Q13 ctx
\textbf{Read the following passage and answer the questions 13--16.}\par\smallskip
Both plants and animals of many sorts show remarkable changes in form, structure, growth habits, and even mode of reproduction in becoming adapted to different climatic environment, types of food supply, or mode of living. This divergence in response to evolution is commonly expressed by altering the form and function of some part or parts of the organism, the original identity of which is clearly discernible. For example, the creeping foot of the snail is seen in related marine pteropods to be modified into a flapping organ useful for swimming, and is changed into prehensile arms that bear suctorial disks in the squids and other cephalopods. The limbs of various mammals are modified according to several different modes of life-for swift running (cursorial) as in the horse and antelope, for swinging in trees (arboreal) as in the monkeys, for digging (fossorial) as in the moles and gophers, for flying (volant) as in the bats, for swimming (aquatic) as in the seals, whales and dolphins, and for other adaptations. The structures or organs that show main change in connection with this adaptive divergence are commonly identified readily as homologous, in spite of great alterations. Thus, the finger and wristbones of a bat and whale, for instance, have virtually nothing in common except that they are definitely equivalent elements of the mammalian limb.

%% Q13
The passage describes how organisms diverge in form in \emph{adapting} to different environments: \textbf{Adaptive Divergence} (the term the passage itself uses).

%% Q14
The passage explains that climate, food supply and mode of living cause change (I), and it relates the limbs of horses (running) and seals (swimming) as homologous (III). It does not explain the theory of evolution (II). \textbf{I and III only}.

%% Q15
\textbf{Homologous} organs are equivalent or \textbf{corresponding} parts that share an origin despite different functions.

%% Q16
The author presents facts and examples without personal bias: an \textbf{objective} style.

%% Q17 ctx
\textbf{Read the following passage and answer the questions 17--22.}\par\smallskip
The atmosphere is a mixture of several gases. There are about ten chemical elements which remain permanently in gaseous form in the atmosphere under all natural conditions. Of these permanent gases, oxygen makes up about 21 percent and nitrogen about 78 percent. Several other gases, such as argon, carbon dioxide, hydrogen, neon, krypton and Xenon, comprise the remaining 1 percent of the volume of dry air. The amount of water vapour, and its variations in amount and distribution, are of extraordinary importance in weather changes. Atmospheric gases hold in suspension great quantities of dust, pollen, smoke and other impurities which are always present in considerable, but variable amounts.

The atmosphere has no definite upper limits but gradually thins until it becomes imperceptible. Until recently it was assumed that the air above the first few miles gradually grew thinner and colder at a constant rate. It was also assumed that upper air had little influence on weather changes. Recent studies of the upper atmosphere, currently being conducted by earth satellites and missile probings, have shown these assumptions to be incorrect. The atmosphere has three well-defined strata.

The layer of the air next to the earth, which extends upward for about 10 miles, is known as troposphere. On the whole, it makes up about 75 percent of all the weight of the atmosphere. It is the warmest part of the atmosphere because most of the solar radiation is absorbed by the earth's surface, which warms the air immediately surrounding it. A steady decrease of temperature with increasing elevation is a most striking characteristic. The upper layers are colder because of their great distance from the earth's surface and rapid radiation of heat into space. The temperatures within the troposphere decrease about 3.5 degrees per 1,000 foot increase in altitude. Within the troposphere, winds and air currents distribute heat and moisture. Strong winds, called jet streams, are located at the upper levels of the troposphere. These jet streams are both complex and wide spread in occurrence. They normally show a wave-shaped pattern and move from west to east at velocities of 150 mph, but velocities as high as 400 mph have been noted. The influences of changing locations and strengths of jet streams upon weather conditions and patterns are no doubt considerable. Current intensive research may eventually reveal their true significance.

Above the troposphere to a height of about 50 miles is a zone called the stratosphere. The stratosphere is separated from the troposphere by a zone of uniform temperatures called tropopause. Within the lower portions of the stratosphere is a layer of ozone gases which filters outmost of the ultraviolet rays from the sun. The ozone layer varies with air pressure. If this zone was not there, the full blast of the sun's ultraviolet light would burn our skin, blind our eyes, and eventually result in our destruction. Within the stratosphere, the temperature and atmospheric composition are relatively uniform.

The layer upward of about 50 miles is the most fascinating but the least known of these three strata. It is called ionosphere because it consists of electrically charged particles called ions, thrown from the sun. The northern lights (aurora borealis) originate within this highly charged portion of the atmosphere. Its effect upon weather conditions, if any, is as yet unknown.

%% Q17
The passage describes the composition and layers of the atmosphere and repeatedly relates them to weather: \textbf{The Atmosphere: its nature and importance to our weather}.

%% Q18
The \textbf{ozone layer} filters the Sun's ultraviolet rays; without it ``the full blast of the sun's ultraviolet light would \ldots\ result in our destruction''.

%% Q19
Only question I is answered (the troposphere and stratosphere are contrasted in temperature and composition). The ionosphere's effect on weather is ``as yet unknown'' (II), and how satellites study the atmosphere is not explained (III). \textbf{I only}.

%% Q20
The troposphere is warmest ``because most of the solar radiation is absorbed by the earth's surface, which warms the air immediately surrounding it'', i.e.\ it is \textbf{warmed by the earth's heat}.

%% Q21
The passage stresses the \emph{variations} in the amount of water vapour, so ``a \textbf{definite} amount of water vapour'' is the exception.

%% Q22
Temperature falls about $3.5^\circ$ per 1,000 ft, so at 12,000 ft: $12\times3.5=\mathbf{42}^\circ$ colder than on the ground.

%% Q23
The speaker uses words in their \emph{opposite} sense (sky is ``down'', winter is ``hot''), so a criminal would be called a \textbf{saint}.

%% Q24
Dolly's mother is John's paternal aunt, so Dolly's father is the husband of John's aunt, i.e.\ John's uncle (by marriage). John is his \textbf{nephew}.

%% Q25
20 : 21 are consecutive numbers, corresponding to T : U, consecutive letters. Similarly \textbf{A : B} are consecutive (positions 1, 2).

%% Q26 ctx
\textbf{Study the pie-chart and answer the questions 26 to 30.} (The percentages below are those implied by the questions and official answers; the original chart was printed without values.)
\begin{center}
\begin{tikzpicture}\pie[radius=1.4,text=pin,font=\scriptsize,color={brand!35,brand!10,brand!25,brand!45,brand!5}]{25/Engineering,12.5/Medical,37.5/Arts,12.5/Commerce,12.5/Science}\end{tikzpicture}
\end{center}

%% Q26
Medical $+$ Engineering $=12.5\%+25\%=\mathbf{37.5\%}$.

%% Q27
Arts $+$ Science $=37.5\%+12.5\%=\mathbf{50\%}$.

%% Q28
Arts $+$ Engineering $=37.5\%+25\%=\mathbf{62.5\%}$.

%% Q29
Engineering $=25\%=320$, so $1\%=12.8$ and Arts $=37.5\%=37.5\times12.8=\mathbf{480}$.

%% Q30
Commerce $+$ Arts $+$ Medical $=12.5+37.5+12.5=62.5\%$ of $1280=\mathbf{800}$.

%% Q31
\begin{center}\begin{tikzpicture}[scale=0.55,font=\scriptsize]\draw (0,0) circle(1.8); \node at (0,1.35){Eagles}; \draw (0,-0.3) circle(1.1); \node at (0,0.4){Watches}; \draw[fill=brand!15] (0,-0.6) circle(0.55) node{Streets};\end{tikzpicture}\end{center}
All streets are watches and all watches are eagles, so \textbf{all streets are eagles} (I follows). II (all watches are streets) is the illicit converse. \textbf{Only I}.

%% Q32
All boys are mothers and all mothers are fathers, so \textbf{all boys are fathers} (II follows). I (all mothers are boys) is an illicit converse. \textbf{Only II}.

%% Q33 note
Strictly, ``Some birds are clouds'' converts validly to ``Some clouds are birds'', so II follows. The official key gives (D).

%% Q33
``Monkey is a bird'' and ``some birds are clouds'' do not show that the monkey is among those birds that are clouds, so I does not follow. II is the simple converse of the second premise. The official key, however, gives \textbf{neither follows}.

%% Q34 ctx
Expected (E) as well as actual (A) increase in export of different items are given in the following table (values in crores). Take into account the proportions of expected and actual export.
\begin{center}\small
\begin{tabular}{l*{8}{c}}\toprule
& \multicolumn{2}{c}{1990} & \multicolumn{2}{c}{1991} & \multicolumn{2}{c}{1992} & \multicolumn{2}{c}{1993}\\\cmidrule(lr){2-3}\cmidrule(lr){4-5}\cmidrule(lr){6-7}\cmidrule(lr){8-9}
\textbf{Items} & E & A & E & A & E & A & E & A\\\midrule
Rice & 90 & 60 & 120 & 60 & 150 & 60 & 200 & 60\\ Sugar & 100 & 70 & 200 & 140 & 300 & 210 & 400 & 280\\ Wheat & 100 & 75 & 200 & 150 & 300 & 200 & 400 & 300\\ Textile & 500 & 450 & 700 & 600 & 800 & 700 & 1000 & 850\\ Machinery & 1000 & 600 & 1500 & 700 & 2000 & 1500 & 2500 & 1500\\ Electronic goods & 1500 & 1200 & 2000 & 1500 & 2500 & 2000 & 3000 & 2000\\\bottomrule
\end{tabular}\\[4pt]
\textbf{Actual as \% of expected:}\\[2pt]
\begin{tabular}{lcccc}\toprule & 1990 & 1991 & 1992 & 1993\\\midrule
Rice & 66.7 & 50.0 & 40.0 & 30.0\\ Sugar & 70.0 & 70.0 & 70.0 & 70.0\\ Wheat & 75.0 & 75.0 & 66.7 & 75.0\\ Textile & 90.0 & 85.7 & 87.5 & 85.0\\ Machinery & 60.0 & 46.7 & 75.0 & 60.0\\ Electronic goods & 80.0 & 75.0 & 80.0 & 66.7\\\bottomrule
\end{tabular}
\end{center}

%% Q34
Only \textbf{rice} falls every year: 66.7\% $\rightarrow$ 50\% $\rightarrow$ 40\% $\rightarrow$ 30\%.

%% Q35
In 1990: Machinery 60\%, \textbf{Textile 90\%}, Wheat 75\%, Sugar 70\%. Textile is highest.

%% Q36
Wheat: 75, 75, \textbf{66.7}, 75. The lowest is \textbf{1992}.

%% Q37
Sugar's percentage is 70\% in every year, so there is no single lowest year: \textbf{none of these}.

%% Q38
1993: \textbf{Rice 30\%}, Sugar 70\%, Wheat 75\%, Machinery 60\%. Rice is lowest.

%% Q39
See, hear and smell are \emph{sense perceptions}; \textbf{think} is a mental (cognitive) process.

%% Q40
TRP, SQO and YWU move backwards in steps of 2; \textbf{TVX} moves forwards, so it is different.

%% Q41
Teaching means \textbf{creating an environment conducive to learning}, not merely giving information or dictating notes.

%% Q42
Per the key, the untrue statement is (D). The lecture method is widely used, fairly structured and efficient for conveying information; statement (D) overstates its standing.

%% Q43
Good teaching is reflected in the \textbf{meaningful questions asked by students}, which show understanding and curiosity.

%% Q44
The main objective of higher education teaching is to \textbf{develop reasoning and thinking abilities}.

%% Q45
Noise in class may come from interaction among students \emph{and/or} between teacher and students: \textbf{(A) and/or (B)}.

%% Q46
Reading a book is communication between the author (sender) and the reader (receiver) through the printed text (channel): \textbf{there is communication} (one-way).

%% Q47
Facial expressions are a \textbf{non-verbal channel} through which students send messages (interest, confusion, boredom) back to the teacher. The key therefore relates them to the \textbf{channel} of communication.

%% Q48
Communication is the \textbf{exchange of information, feelings, emotions and views}.

%% Q49
Indian classrooms are typically dominated by the lecture: \textbf{one-way verbal communication}.

%% Q50
Two people engrossed in their own thoughts are not exchanging any message, so there is \textbf{no communication} between them.
